\section{Current reading of the delivered residue--constituent note}

The complete delivered note, equations (RC1)--(RC53), is retained
unchanged. The following statements accompany that historical source
in the current integrated edition.

First, the translation source-byte comparison has already been resolved.
The dated translation seal proves that the reviewed source becomes the
final source by the single geometry change from a \(25\)-millimetre
margin to a \(23\)-millimetre margin. Reversing that exact byte-string
substitution reproduces the reviewed SHA-256 value. The later README
records the completed delivery. The historical final paragraph of the
delivered note therefore records an intake question that the existing
seal answers.

Second, the precise derivative convention in (RC28) is the real
axis \(t=x\in\mathbb R\), with the displayed real \(u>0\).
For \(Y=\Pi I\) and \(H_F=Y^*Y\), the period equation gives
\[
 Y'(0)=-Y(0)A_F/u,\qquad
 \left.\frac{d}{dx}H_F(x)\right|_0
       =-(A_F^*H_F(0)+H_F(0)A_F)/u.
\]
Taking its logarithmic determinant gives
\(-2\operatorname{Re}\operatorname{Tr}A_F/u\), precisely (RC28).
The corresponding Wirtinger derivative is
\(-\operatorname{Tr}A_F/u\). The original fixed denominator
\(\det G_{N,F}\) contributes zero to both derivatives.

Third, (RC27) has the correct mixed fourth derivative and factorial.
Its exact relation to real-axis differentiation is (RCX26)--(RCX32):
\[
 \psi_{tt\bar t\bar t}(0)=\mathfrak c_F(H(0))/|u|^2,\qquad
 \left.\frac{d^4}{dx^4}\psi(x)\right|_0
   =2\operatorname{Re}\phi^{(4)}(0)
      +6\mathfrak c_F(H(0))/|u|^2 .
\]
Here \(\phi(t)=\log\det(I^*\Pi(0)^*\Pi(t)I)\), and its full
fourth derivative is evaluated by the displayed noncommutative
period recursion in the companion. The coefficient of
\(t^2\bar t^2\) is the mixed derivative divided by \(2!2!=4\).

Fourth, the source equation (RC41) is evaluated on the admitted
strong-Schwartz tensor domain stable under the original
\(D^{(k)}\) and its square. The exact lift
\(\widehat R=r_NRj_E\) satisfies
\(j_E\widehat R=Rj_E\) and \(\widehat R d=0\).
The complete primitive relating the quadratic actions is
\[
 ((D^{(k)}+t\widehat R)^2-k(D^{(k)}+t\widehat R))r_N
     -r_N(A_t^2-kA_t)
       =d(D_0K+KA_t-kK),
\]
where \(D^{(k)}r_N-r_NA=dK\) and \(D^{(k)}d=dD_0\).
Every term has the original source and primitive target.

The shared general rank and curvature proofs were already integrated
as (SC9)--(SC25); the source comparison in (RC32)--(RC33) is
also (LC14)--(LC15). The new companion adds the universal
endomorphism inverse, the exact metric differential and sharpened
endpoint estimate, the complete higher period jets, and the forced
Laplacian map. The delivered twelve-method suite has been executed
in normal and optimized modes with six substantive negative-control
failures. Its translation method checks the full \(q\)-dimensional
trace shift; the constituent \(p\)-dimensional formula is proved by
the full frame transport in (RCX36)--(RCX39).
